% ----- CHAUPAI 2 -----

\newpage

\subsection*{\deva{द्वितीय चौपाई} (Second Chaupai)}
\addcontentsline{toc}{subsection}{\deva{द्वितीय चौपाई} (Second Chaupai) — \deva{राम दूत...}}

\begin{minipage}[t]{0.55\textwidth}
\vspace{0pt}

{\large \linespread{1.5}\selectfont
\deva{राम दूत अतुलित बल धामा।}\\
\deva{अंजनि-पुत्र पवन सुत नामा॥}
\par}

\vspace{0.5em}

{\large \linespread{1.5}\selectfont
\telu{రామ దూత అతులిత బల ధామా।}\\
\telu{అంజని పుత్ర పవన సుత నామా॥}
\par}

\vspace{0.5em}

{\large \linespread{1.3}\selectfont
\textit{rāma dūta atulita bala dhāmā |}\\
\textit{añjani-putra pavana suta nāmā ||}
\par}
\end{minipage}%
\hfill
\begin{minipage}[t]{0.42\textwidth}
\vspace{0pt}
\begin{tcolorbox}[anchorbox]
\textbf{Functional Leadership:} True leadership is about capability and trust, not titles. When millions of soldiers were dispatched to find Sita, Prince Angada was the official commander of the party, yet Lord Rama handed his personal signet ring to Hanuman alone. Rama recognized Hanuman as the functional leader whose competence, intellect, and empathy would guide the team to success.
\vspace{0.5em}\newline Out of millions, Lord Rama entrusted his signet ring uniquely to Hanuman because of his unmatched leadership and reliability.\newline\null\hfill\href{https://www.valmiki.iitk.ac.in/sloka?field_kanda_tid=4&language=dv&field_sarga_value=44}{\textit{Kishkindha Kanda, 44:11-13}}
\end{tcolorbox}
\end{minipage}

\vspace{0.5em}
\subsubsection*{\textbf{Visual Rhythm Map:}}
\begin{table}[h!]
\centering
\renewcommand{\arraystretch}{1.2}
\setlength{\tabcolsep}{0pt}
\begin{tabularx}{\textwidth}{|*{16}{>{\centering\arraybackslash\hsize=1.0\hsize}X|}}
\hline
\multicolumn{3}{|>{\centering\arraybackslash\hsize=3.0\hsize}X|}{\textbf{\guru{रा}\laghu{म}} \par (3)} &
\multicolumn{3}{>{\centering\arraybackslash\hsize=3.0\hsize}X|}{\textbf{\guru{दू}\laghu{त}} \par (3)} &
\multicolumn{4}{>{\centering\arraybackslash\hsize=4.0\hsize}X|}{\textbf{\laghu{अ}\laghu{तु}\laghu{लि}\laghu{त}} \par (4)} &
\multicolumn{2}{>{\centering\arraybackslash\hsize=2.0\hsize}X|}{\textbf{\laghu{ब}\laghu{ल}} \par (2)} &
\multicolumn{4}{>{\centering\arraybackslash\hsize=4.0\hsize}X|}{\textbf{\guru{धा}\guru{मा}} \par (4)} \\
\hline
\end{tabularx}
\vspace{-1.5em}
\end{table}

\begin{table}[h!]
\centering
\renewcommand{\arraystretch}{1.2}
\setlength{\tabcolsep}{0pt}
\begin{tabularx}{\textwidth}{|*{16}{>{\centering\arraybackslash\hsize=1.0\hsize}X|}}
\hline
\multicolumn{4}{|>{\centering\arraybackslash\hsize=4.0\hsize}X|}{\textbf{\guru{अं}\laghu{ज}\laghu{नि}} \par (4)} &
\multicolumn{3}{>{\centering\arraybackslash\hsize=3.0\hsize}X|}{\textbf{\guru{पु}\laghu{त्र}} \par (3)} &
\multicolumn{3}{>{\centering\arraybackslash\hsize=3.0\hsize}X|}{\textbf{\laghu{प}\laghu{व}\laghu{न}} \par (3)} &
\multicolumn{2}{>{\centering\arraybackslash\hsize=2.0\hsize}X|}{\textbf{\laghu{सु}\laghu{त}} \par (2)} &
\multicolumn{4}{>{\centering\arraybackslash\hsize=4.0\hsize}X|}{\textbf{\guru{ना}\guru{मा}} \par (4)} \\
\hline
\end{tabularx}
\vspace{-1.5em}
\end{table}

\vspace{1em}

    \textbf{ \deva{सरल-संस्कृतम्}:}\\
 \deva{हे रामदूत! त्वम् अतुलित-बलस्य धाम (आश्रयः) असि।}\\
 \deva{हे अञ्जनिपुत्र! त्वं पवनसुत-नाम्ना प्रसिद्धः असि॥}\\

O Emissary of Rama! You are the abode of incomparable strength.\\
O Son of Anjani! You are famous by the name 'Son of the Wind'.


\vspace{1em}
\newpage
\textbf{\deva{प्रतिपदार्थः} (Meaning of each word):}
\begin{table}[h!]
\centering
\renewcommand{\arraystretch}{1.2}
\begin{tabularx}{\textwidth}{@{} l l X X @{}}
\toprule
\textbf{Awadhi} & \textbf{Sanskrit} & \textbf{English} & \textbf{Grammar \& Notes} \\
\midrule
\deva{राम दूत} / \telu{రామ దూత} / \textit{rāma dūta}  &  \deva{श्रीरामस्य दूतः}  &  Emissary of Rama  &   \\
\deva{अतुलित} / \telu{అతులిత} / \textit{atulita}  &  \deva{अतुलितम्}  &  Incomparable  &   \\
\deva{बल} / \telu{బల} / \textit{bala}  &  \deva{बलम्}  &  Strength  &   \\
\deva{धामा} / \telu{ధామా} / \textit{dhāmā}  &  \deva{धाम / आश्रयः}  &  Abode / Repository  & Meter: Added 'aa' for rhyme \\
\deva{अंजनि पुत्र} / \telu{అంజని పుత్ర} / \textit{añjani putra}  &  \deva{अञ्जनायाः पुत्रः}  &  Son of Anjani  &   \\
\deva{पवन सुत} / \telu{పవన సుత} / \textit{pavana suta}  &  \deva{पवनस्य सुतः}  &  Son of the Wind  &   \\
\deva{नामा} / \telu{నామా} / \textit{nāmā}  &  \deva{नाम्ना}  &  By the name  & Meter: Added 'aa' for rhyme \\
\bottomrule
\end{tabularx}
\end{table}

\begin{tcolorbox}[poetrybox]
\textbf{Praasa (The Rhyme Engine):}\\
Note the elongated \deva{आ} (aa) at the end of \deva{धामा} (\textit{dhāmā}) and \deva{नामा} (\textit{nāmā}). Sanskrit/Hindi would use \deva{धाम} (Dhaam) and \deva{नाम} (Naam), but Awadhi uses the heavy 2-beat suffix to complete the 16-beat line and create the musical bounce.
\end{tcolorbox}

\begin{tcolorbox}[gotchabox]
\begin{itemize}
    \item \textbf{Conjunct Consonant Weight (\deva{पुत्र}):} The short 'u' in \deva{पु} gets weighed down and becomes a 2-beat Guru because it crashes into the conjunct consonant \deva{त्र} (t+ra).
\end{itemize}
\end{tcolorbox}



\newpage
